\PBchapter{O Tempo}
% \PByear{1977}
Para falar de uma jornada, precisamos falar sobre o Tempo.

Existem literaturas, artigos e pesquisas que costumam classificar as pessoas por gerações.

Para mim, isso sempre pareceu algo muito mais estadunidense — até porque essa lógica nasce no pós-guerra. 

Então, dentro dessa nomenclatura, temos os Baby Boomers (1946–1964), a Geração X (1965–1980), os Millennials ou Geração Y (1981–1996), e por aí vai.

Seguindo essa linha, eu seria da Geração X. Afinal, nasci em Abril de 1977.

Mas… e o Brasil nisso tudo?

Meus avós nasceram entre as décadas de 10 e 30, como ficam nessa escala?

Meus pais nasceram nos anos 50, loga, são da geração Baby Boomers?

Como essas experiências que tiveram em suas épocas, com seus desafios, afetaram direta ou indiretamente a minha vida?

Eu realmente sou da Geração X?

Difícil responder.

Mas dá para refletir.

Meus avós maternos nasceram na Bahia.

%\PByear{1910}

Meu avô nasceu por volta de 1925. 

Ninguém sabe ao certo o ano. Tanto que, no registro, consta 1º de janeiro — uma data “padrão” para quem não tinha informação precisa. Nasceu em Mucugê, interior da Bahia.


Minha avó nasceu em 1933, mas foi registrada como 1930 — já em São Paulo. 

Ela nasceu em um lugar chamado Jurema na Bahia, que nunca consegui localizar exatamente. Provavelmente foi um distrito, povoado ou até uma fazenda.

E sim, meus avós são fruto da mistura brasileira: brancos, negros e povos originários.

A trajetória deles, imagino, não foge muito da de tantas outras pessoas que migraram do Nordeste.

No caso do meu avô, ele saiu de Mucugê porque alguém "descobriu" diamantes na região.

E isso bastou.

Virou um caos.

Em algum momento, meus avós se encontraram em São Paulo — não sei dizer exatamente quando, nem onde.

Mas sei que, depois disso, meu avô foi trabalhar na Companhia Melhoramentos Norte do Paraná, conhecida também como Paraná Plantations.

Sim, derrubando árvores.

Isso já no Noroeste do Paraná, mais especificamente na região de Umuarama.

Estamos falando de algo entre 1953 e 1973.

Tenho essa referência porque minha mãe nasceu em 1953, na cidade de Nova Esperança, no norte do Paraná. 

Em 1973, a família se mudou para Cascavel, no oeste do estado do Paraná.

Agora, do lado paterno, a história é mais fragmentada.

Consegui montar uma espécie de prévia da árvore genealógica, mas com muitas lacunas.

E aqui vale um parêntese.

Meus pais se separaram quando eu tinha oito anos. Tive pouco contato com a família paterna. Na verdade, nem meu pai teve muito.

Segundo ele contou à minha mãe, fugiu de casa aos quatorze anos.

Fecha parêntese.

O que consegui levantar foi o seguinte.

Meus avós paternos se casaram por volta de 1943, em uma região hoje conhecida como Alto Uruguai, distrito de Tiradentes do Sul, no oeste do Rio Grande do Sul.

Ambos nasceram por volta de 1910.

Na ocasião do casamento, aproveitaram para registrar os filhos que já tinham — algo comum na época. 

%\PByear{1950}

Meu pai, que nasceu em 1950, nem aparece nesse registro.

Eu sei… pode parecer que estou me perdendo.

Que nada se conecta.

Mas continua comigo.

Porque essa história — essa micro-história — é justamente o contexto das reflexões que me trouxeram até a tal parede branca.

Se usamos o fim da Segunda Guerra Mundial como marco para definir gerações no mundo, por que não olhar para marcos históricos do Brasil?

Pode parecer viagem.

Mas nós tivemos a nossa guerra.

Entre 1864 e 1870, ocorreu a Guerra do Paraguai — também chamada de Guerra da Tríplice Aliança.

Talvez, para muita gente, isso não tenha impacto direto.

Mas, nas minhas reflexões, têm.

Após o fim da guerra, o Império Brasileiro passou a incentivar a ocupação das regiões de fronteira.

E adivinha?

A região do Alto Uruguai, onde minha família paterna aparece, está exatamente nessa faixa de fronteira — próxima à Argentina e, historicamente, ligada a disputas territoriais com o Paraguai.

Não tenho como afirmar com precisão, mas tudo indica um movimento migratório nesse contexto.

O registro mais antigo que encontrei é o de uma trisavó, nascida por volta de 1885.

Depois disso, o que existe é mais suposição do que certeza.

Já do lado materno, a história passa por outro tipo de movimento.

E aqui entra a Paraná Plantations — que eu não mencionei por acaso.

Essa empresa, de origem inglesa, comprou grandes extensões de terras e foi responsável pela colonização do Norte do Paraná. Cidades como Londrina, Maringá, Apucarana, Cambé, Cianorte e tantas outras surgiram a partir desse processo.

Em 1942, durante a Segunda Guerra Mundial, a empresa foi nacionalizada.

Ou seja, a Guerra Mundial também atravessa essa história.

E, nesse cenário, meu avô acaba indo parar em Umuarama.

Durante minhas pesquisas sobre essa região e esse período, uma coisa me chamou a atenção: a quantidade de registros de mortes violentas — brigas, esfaqueamentos, conflitos.

Não necessariamente da minha família.

Mas o fato de muitos registros nem indicarem a origem dos que morreram, diz muito sobre o tipo de ambiente da época.

Era um território em formação.

Sem raízes.

Sem estabilidade.

Chegamos então aos meus pais.

Meu pai sai da região de Três Passos, no Rio Grande do Sul, e vai para Planalto, cidade do sudoeste do Paraná, divisa com Santa Catarina e Argentina.

Ali aprende a profissão de mecânico.

Estamos falando algo entre 1964 e 1975.

Não sei muito sobre a história dele nesse período.

Já minha mãe chega a Cascavel em 1973, com 20 anos. 

Costureira. 

E como era comum na época não pode estudar e nem trabalhar.

Eles se conhecem.

%\PByear{1970}

Se casam em 1975.

No mesmo ano, nasceu minha irmã — de sete meses. E sim, na época achavam que meus pais “se adiantaram” (sic).

Nasci em 1977.

E é aqui que o título deste capítulo começa a fazer sentido.

A gente costuma pensar que a vida começa no nosso nascimento.

Mas não começa.

Existe todo um tempo antes.

Um acúmulo de decisões, movimentos, acasos, guerras, migrações — que, se fossem diferentes, talvez nem permitissem que a gente existisse.

Então volto à pergunta.

De que geração eu sou?

Sinceramente, não sei.

Talvez mais importante do que isso seja olhar para esse recorte de tempo — de 1870, fim da Guerra do Paraguai, até 1977.

São pouco mais de cem anos.

E, nesse intervalo, o Brasil saiu de um Império, passou por crises, golpes, mudanças profundas… e chegou a uma ditadura militar.

Mudanças colossais.

O quanto isso impactou os lugares por onde minha família passou?

O quanto isso moldou as pessoas?

Eu sei que, na minha família, isso deixou marcas.

Na falta de vínculos.

Na dificuldade de afeto.

Na ausência de pertencimento.

Não sei se eles percebem isso.

Mas vivem isso.

E isso aparece na forma como nos encontramos — quase sempre em momentos de dor, perda ou necessidade.
Isso também chegou até mim.

E, voltando à ideia de geração…

Existe uma descrição comum da Geração X: pessoas equilibradas, felizes, com boa relação entre vida pessoal e profissional, produtivas, empreendedoras.

Eu não me reconheço nisso.

Não tive esse equilíbrio.

Nem antes, nem agora, em 2026.

Não sou empreendedor.

Talvez a única coisa que se mantenha seja a produtividade — mas construída à base de medo, ansiedade e pressão.

E isso molda a forma como eu vejo o mundo.

Ou talvez explique por que, às vezes, eu sinto que não vejo com clareza.

Como se existisse um descompasso.

Como se o tempo, ao invés de organizar, desorganiza, bagunça, complica.

%\PBreflection{A gente costuma pensar que a vida começa no nosso nascimento. Mas ela começa muito antes de nós.}
